problem:  
**Problem (Rigidity of locally homogeneous non‑parallel Spin(7) structures).**  
Let \((M^8, g, \Phi)\) be a connected oriented Riemannian 8‑manifold carrying a Spin(7)‑structure given by a positive 4‑form \(\Phi\).  
For a general Spin(7)‑structure, the intrinsic torsion is encoded by a 1‑form \(\theta\) (the **Lee form**) and a component \(\tau_{48}\) taking values in an irreducible 48‑dimensional Spin(7) representation, so that  
\[
d\Phi = \theta\wedge\Phi + \tau_{48}.
\]  
The structure is **torsion‑free** (and thus holonomy contained in Spin(7)) precisely when both \(\theta=0\) and \(\tau_{48}=0\).

Assume \((M,g,\Phi)\) is **locally homogeneous**, i.e., modeled on the left‑invariant geometry of some 8‑dimensional Lie group with a left‑invariant positive 4‑form \(\Phi\).

**Open problem.**  
Determine whether there exist **non‑flat, locally homogeneous** Spin(7)‑structures with **exact Lee form** (\(\theta = d f\) for some local function \(f\)) and \(\tau_{48}=0\) (“conformally parallel” Spin(7)-structures) that are Einstein but not locally symmetric.  
Equivalently:  
- If a locally homogeneous Spin(7)-structure is conformally parallel, must it be locally conformally flat or locally symmetric?  
- Can a nontrivial conformal factor yield an Einstein, locally homogeneous but non‑flat metric compatible with Spin(7) geometry?

---

Why is it a "good" problem:  
1. **Mathematically sound and precise.**  
   The problem now uses the correct and standard torsion decomposition for Spin(7)‑structures—via the Lee form and the \(\tau_{48}\)‑component—ensuring logical validity and grounding in established differential‑geometric theory.

2. **Unifies local homogeneity and conformal Spin(7) theory.**  
   The conformally parallel case (\(\tau_{48}=0\)) leads to a first‑order nonlinear PDE linking the conformal factor and the Spin(7) 4‑form. Combining this with the algebraic curvature constraints coming from local homogeneity produces a rich system intertwining Lie group geometry, conformal structure, and exceptional holonomy.

3. **Genuinely open territory.**  
   While torsion‑free (Ricci‑flat) and fully parallel Spin(7)‑structures are known to be rigid (no nonflat homogeneous examples), virtually nothing is known about the **locally homogeneous conformally parallel** class. Even the existence of algebraic Lie‑algebra models satisfying the combined Einstein and conformal Spin(7) equations is unresolved.

4. **Cross‑disciplinary depth.**  
   The problem naturally connects:
   - representation‑theoretic analysis (\(\tau_{48}\) and Spin(7) decomposition of forms),
   - conformal differential geometry (Lee form and local rescaling behavior),
   - Lie algebraic curvature classification (Ambrose–Singer framework for local homogeneity),
   - and Ricci curvature analysis under Spin(7) constraints.  
   It thereby bridges major subfields of modern Riemannian geometry.

5. **Potentially transformative outcomes.**  
   - A rigidity proof would extend the known flatness results for homogeneous exceptional holonomy much further, showing that *even up to conformal change*, local homogeneity annihilates nontrivial Spin(7) holonomy.  
   - A counterexample would yield the **first known explicit conformally parallel Spin(7)** homogeneous model, unlocking a new class of Einstein metrics and impacting calibrations and compactification models in physics.

6. **Simplicity masking depth.**  
   The core question—*Can a locally homogeneous, conformally parallel Spin(7) geometry be nontrivially Einstein?*—is succinct and elementary to state, yet resolving it demands harmonizing algebraic and analytic tools across special holonomy, conformal, and homogeneous geometries.  
   This blend of brevity, internal consistency, cross‑field connectivity, and high‑stakes insight exemplifies a **“narrow entrance, profound depth”** research‑level mathematical problem.